\section{Pendahuluan}
Arsitektur microservice memecah sebuah aplikasi menjadi layanan (service) yang di-deploy secara mandiri dan saling memanggil. Pola ini meningkatkan skalabilitas dan kemudahan pemeliharaan, tetapi anomali pada satu service dapat merambat ke service yang bergantung padanya dan akhirnya menyebabkan kegagalan sistem secara keseluruhan \cite{lin2024run}. Ketika insiden terjadi, insinyur on-call harus menemukan service yang menjadi sumber gangguan. Tugas ini disebut Root Cause Localization (RCL), yaitu bagian dari Root Cause Analysis (RCA) \cite{sun2024art}. Penelusuran manual berjalan lambat, dan banyaknya metrik, log, serta trace yang dihasilkan sistem yang sedang berjalan membuat dukungan otomatis diperlukan \cite{pham2024rcaeval}.

Metode RCL terawasi bergantung pada insiden historis berlabel. Label tersebut langka dan mahal karena akar masalah harus dikonfirmasi oleh operator berpengalaman \cite{sun2024art}. Self-Supervised Learning (SSL) menawarkan cara memanfaatkan data pemantauan tanpa label: encoder dilatih lebih dulu dengan tugas pretext, kemudian diadaptasi dengan sedikit data berlabel. Metode SSL deret waktu terkini, seperti TS2Vec yang berbasis kontrastif \cite{yue2022ts2vec} dan TimeMAE yang berbasis masked autoencoder \cite{cheng2026timemae}, menargetkan representasi yang dapat ditransfer. TimeMAE melaporkan keunggulan pada kondisi label terbatas dan ketika encoder yang di-pre-train pada satu dataset di-fine-tune pada dataset lain, sedangkan TS2Vec menerapkan encoder yang dilatih pada satu dataset untuk deteksi anomali pada dua dataset lain. Namun, keduanya dievaluasi pada tolok ukur deret waktu umum, yaitu sebuah sampel berupa satu deret dengan label kelas, target peramalan, atau penanda anomali. RCL berbeda: satu kasus memuat banyak service yang saling bergantung, dan keluarannya adalah pilihan di antara service tersebut.

Metode analisis akar masalah RUN \cite{lin2024run} dan ART \cite{sun2024art} memodelkan lingkungan microservice secara langsung. RUN memakai backbone peramalan yang di-pre-train secara kontrastif di dalam pipeline causal discovery per kasus, dan dievaluasi pada data sintetis serta satu test bed. ART mempelajari perilaku sistem dari data normal dan dievaluasi pada dua dataset, salah satunya tidak publik. Tidak satu pun dari kedua studi itu menguji apakah apa yang dipelajari pada satu dataset RCA bermanfaat pada dataset lain.

\paragraph{Celah penelitian.} Sepengetahuan kami, apakah representasi self-supervised yang di-pre-train pada satu dataset RCA microservice dapat ditransfer ke dataset lain belum pernah diteliti. Bukti yang ada terpisah: encoder SSL divalidasi pada tolok ukur deret waktu umum, sedangkan metode RCA divalidasi satu dataset demi satu dataset. Studi transfer pada tolok ukur bersama juga membutuhkan kontrol yang sering tidak ada, yaitu encoder yang diinisialisasi acak dengan arsitektur sama untuk memisahkan manfaat pre-training dari manfaat fine-tuning terawasi, serta baseline heuristik sederhana.

\paragraph{Tujuan.} Penelitian ini mengkaji transfer representasi deret waktu self-supervised untuk lokalisasi root-cause service pada tolok ukur RCAEval \cite{pham2024rcaeval}. Tiga pendekatan representasi dibandingkan: TS2Vec (kontrastif), TimeMAE (masked autoencoding), dan adaptasi tahap pre-training RUN. RUN dalam bentuk aslinya dievaluasi sebagai baseline RCA tanpa transfer. Karena mekanisme pelatihan ketiga metode berbeda, analisis memisahkan tiga pertanyaan, dan satu protokol tunggal tidak dipaksakan kepada semua metode.

\paragraph{Pertanyaan penelitian.}
\begin{description}
\item[RQ1] Dapatkah encoder yang di-pre-train dengan SSL pada satu dataset RCAEval ditransfer ke dataset lain untuk lokalisasi root-cause service, dibandingkan dengan encoder yang diinisialisasi acak?
\item[RQ2] Bagaimana perbandingan TS2Vec, TimeMAE, dan RUN dari sisi (a) akurasi lokalisasi dan (b) manfaat yang diperoleh dari pre-training dan transfer?
\item[RQ3] Bagaimana proporsi data latih berlabel (10\%, 25\%, 50\%, 100\%) dan strategi fine-tuning (FineLast dibandingkan FineAll) memengaruhi akurasi serta manfaat pre-training?
\end{description}

\paragraph{Kontribusi.} (i) Protokol pembelajaran transfer untuk RCL tingkat service pada RCAEval dengan skenario in-domain dan lintas dataset, kontrol encoder acak dan heuristik, serta uji bootstrap berpasangan. (ii) Eksperimen kontrol yang mengganti data pasca-injeksi dengan data pra-injeksi untuk memeriksa bahwa model membaca anomali, bukan menghafal identitas service. (iii) Perbandingan tiga keluarga SSL dengan aturan yang memisahkan perbandingan metode, efek pre-training, dan efek label serta fine-tuning. (iv) Uraian tentang kondisi ketika pre-training dan transfer tidak membantu. \todo{Konfirmasi kontribusi setelah eksperimen TimeMAE dan RUN selesai; temuan utama ditambahkan ke paragraf ini hanya setelah tersedia.}

Sisa artikel disusun sebagai berikut. Bagian~\ref{sec:related} meninjau penelitian terkait, Bagian~\ref{sec:method} memaparkan metodologi, Bagian~\ref{sec:setup} dataset dan pengaturan eksperimen, Bagian~\ref{sec:results} dan~\ref{sec:discussion} hasil dan pembahasan, dan Bagian~\ref{sec:conclusion} kesimpulan.
